\documentclass[12pt,oneside]{book}

% -------------------------------------------------
% PAQUETES
% -------------------------------------------------
\usepackage[spanish,es-nodecimaldot]{babel}
\usepackage[utf8]{inputenc}
\usepackage[T1]{fontenc}
\usepackage{lmodern}
\usepackage[a4paper,margin=2.5cm]{geometry}
\usepackage{microtype}
\usepackage{setspace}
\usepackage{graphicx}
\usepackage{xcolor}
\usepackage{fancyhdr}
\usepackage{titlesec}
\usepackage{hyperref}
\usepackage{bookmark}
\usepackage{tocloft}
\usepackage{tcolorbox}
\usepackage{amsmath,amssymb}
\usepackage{physics}
\usepackage{enumitem}
\usepackage{csquotes}
\usepackage{listings}

% -------------------------------------------------
% CONFIGURACIÓN VISUAL
% -------------------------------------------------
\definecolor{azulUdeA}{HTML}{026937}
\definecolor{grisClaro}{HTML}{F4F6F8}
\definecolor{grisTexto}{HTML}{2E2E2E}
\definecolor{grisSec}{HTML}{6C757D}

\hypersetup{
    colorlinks=true,
    linkcolor=azulUdeA,
    urlcolor=azulUdeA,
    citecolor=azulUdeA,
    pdftitle={Notas de Clase - Mecánica Celeste},
    pdfauthor={Daniel Soto Villada}
}

\setstretch{1.12}
\setlength{\parskip}{0.4em}
\setlength{\parindent}{0pt}
\setcounter{secnumdepth}{3}
\setcounter{tocdepth}{2}
\setlength{\headheight}{26pt}

% Estilo de capítulos y secciones
\titleformat{\chapter}[display]
  {\bfseries\Huge\color{azulUdeA}}
  {\filleft\Large\chaptername\ \thechapter}
  {1ex}
  {\titlerule\vspace{1ex}\filright}
  [\vspace{1ex}\titlerule]

\titleformat{\section}
  {\bfseries\Large\color{azulUdeA}}
  {\thesection}{0.7em}{}

\titleformat{\subsection}
  {\bfseries\large\color{grisTexto}}
  {\thesubsection}{0.6em}{}

% Encabezados y pies
\pagestyle{fancy}
\fancyhf{}
\fancyhead[L]{\textcolor{grisSec}{\nouppercase{\leftmark}}}
\fancyhead[R]{\textcolor{grisSec}{Daniel Soto Villada}}
\fancyfoot[C]{\textcolor{grisSec}{\thepage}}
\renewcommand{\headrulewidth}{0.4pt}
\renewcommand{\headrule}{\hbox to\headwidth{\color{azulUdeA}\leaders\hrule height \headrulewidth\hfill}}

% Índice más limpio
\renewcommand{\cftchapfont}{\bfseries\color{azulUdeA}}
\renewcommand{\cftsecfont}{\color{grisTexto}}
\renewcommand{\cftchappagefont}{\bfseries\color{azulUdeA}}
\renewcommand{\cftsecpagefont}{\color{grisTexto}}

% Caja para definiciones/ideas importantes
\newtcolorbox{importante}{
  colback=grisClaro,
  colframe=azulUdeA,
  boxrule=0.8pt,
  arc=2mm,
  left=2mm,right=2mm,top=1mm,bottom=1mm
}

% -------------------------------------------------
% DATOS DEL DOCUMENTO
% -------------------------------------------------
\newcommand{\NombreEstudiante}{Daniel Soto Villada}
\newcommand{\Programa}{Estudiante de astronomía}
\newcommand{\Universidad}{Universidad de Antioquia}
\newcommand{\Facultad}{Facultad de Ciencias Exactas y Naturales}
\newcommand{\Instituto}{Instituto de Física}
\newcommand{\TituloDocumento}{Notas de Clase - Mecánica Celeste}

% -------------------------------------------------
% DOCUMENTO
% -------------------------------------------------
\begin{document}
\hypersetup{pageanchor=false}

% ---------- PORTADA ----------
\begin{titlepage}
    \centering
    \vspace*{1.8cm}

    {\Large\textcolor{grisSec}{\Universidad}\par}
    \vspace{0.25cm}
    {\large\textcolor{grisSec}{\Facultad}\par}
    {\large\textcolor{grisSec}{\Instituto}\par}

    \vspace{2.2cm}
    {\Huge\bfseries\textcolor{azulUdeA}{\TituloDocumento}\par}
    \vspace{0.5cm}
    {\Large\textcolor{grisSec}{Plantilla académica profesional}\par}

    \vspace{2.4cm}
    \begin{tcolorbox}[colback=grisClaro,colframe=azulUdeA,boxrule=0.9pt,arc=2mm,width=0.78\textwidth]
        \centering
        {\Large\bfseries \NombreEstudiante}\par
        \vspace{0.2cm}
        {\large \Programa}
    \end{tcolorbox}

    \vfill
    {\large\textcolor{grisSec}{\today}\par}
    \vspace{0.5cm}
\end{titlepage}

% ---------- PÁGINAS PRELIMINARES ----------
\frontmatter
\tableofcontents
\clearpage

% ---------- CONTENIDO ----------
\clearpage
\hypersetup{pageanchor=true}
\mainmatter
\chapter{Introducción}

\section{Objetivo del curso}
El objetivo de estas notas es registrar de forma detallada y rigurosa los conceptos, deducciones y algoritmos vistos en el curso de Mecánica Celeste (2026-1), con énfasis en la comprensión profunda de los fundamentos físicos y matemáticos. En esta clase se aborda el problema relativo de dos cuerpos, pieza clave para reducir el problema de \(N\) cuerpos a un sistema tratable.

\section{Convenciones de apuntes}
\begin{itemize}[leftmargin=2em]
    \item Definiciones clave se destacan en cajas \texttt{importante}.
    \item Teoremas o resultados importantes se derivan paso a paso.
    \item Ejemplos y ejercicios se incluyen para reforzar la comprensión.
\end{itemize}

\begin{importante}
Tip: Usa estas notas para repasar las deducciones analíticas del problema relativo, que permite pasar de 12 variables (6 por cuerpo) a solo 6 variables relativas más el movimiento uniforme del centro de masa.
\end{importante}


\chapter{Clase 13: Sistemas Jerárquicos (parte 1)}

Esta clase pertenece a la Unidad 2 (\emph{El problema de 2-Cuerpos}) y constituye la introducción a los sistemas jerárquicos de \(N\) cuerpos. Se parte de la motivación observacional y teórica de por qué los sistemas reales de \(N\) cuerpos tienden a autoorganizarse en estructuras jerárquicas, y se desarrolla la herramienta matemática (coordenadas de Jacobi) que permite modelarlos eficientemente.

\section{Motivación: autoorganización de los sistemas de \texorpdfstring{$N$}{N} cuerpos}
El problema de \(N\) cuerpos es muy general, pero los sistemas reales tienden a ``autoorganizarse'' después de un tiempo prolongado de interacción.

\begin{itemize}
    \item Con colisiones: se forman cuerpos grandes que crecen \emph{oligárquicamente} (los más masivos crecen más rápido).
    \item Sin colisiones: las interacciones por pares expulsan cuerpos y el sistema se relaja hacia una configuración más ordenada.
\end{itemize}

Las características típicas de estos sistemas autoorganizados son:
\begin{enumerate}
    \item Las masas de los cuerpos difieren mucho entre sí (hay cuerpos muy masivos comparados con otros).
    \item Los cuerpos tienen masas parecidas, pero las distancias entre ellos son muy diferentes (sistemas jerárquicos).
\end{enumerate}

\begin{importante}
En el marco del centro de masa, los sistemas jerárquicos de \(N\) cuerpos evidencian una clara organización por pares. Aunque en otros sistemas de referencia parezcan caóticos, en el CM se observa una estructura ordenada de órbitas anidadas.
\end{importante}

\section{Organización de los sistemas jerárquicos por pares}
Un sistema jerárquico de \(N\) cuerpos puede modelarse como \textbf{\(N-1\) sistemas de dos cuerpos}.

Cada par forma un subsistema cuya órbita relativa se describe con el problema de dos cuerpos ya estudiado. El árbol jerárquico representa la estructura de contención: el sistema total se divide en subsistemas internos y externos.

Ejemplos astrofísicos reales:
\begin{itemize}
    \item Sistema estelar múltiple \textbf{$\alpha$ Gem} (Castor): seis estrellas organizadas en tres pares binarios jerárquicos.
    \item Sistema \textbf{30 Ari}: doble estrella con subcomponentes.
    \item Sistemas planetarios y exoplanetas: un planeta orbita una estrella que a su vez orbita otra estrella lejana.
\end{itemize}

\begin{importante}
En un sistema jerárquico típico, las órbitas internas son mucho más cerradas (y rápidas) que las órbitas externas. Esto permite tratar cada nivel como un problema de dos cuerpos independiente (aproximación jerárquica).
\end{importante}

\section{Tipos de sistemas jerárquicos de dos cuerpos}
Se distinguen dos configuraciones básicas:
\begin{itemize}
    \item \textbf{Sistema central}: un cuerpo masivo en el centro con uno o más cuerpos orbitando a su alrededor (ej. estrella + planeta).
    \item \textbf{Sistema jerárquico propiamente dicho}: dos subsistemas (cada uno con su centro de masa) orbitando mutuamente.
\end{itemize}

\section{Coordenadas de Jacobi}
Para describir eficientemente un sistema jerárquico se utilizan las \textbf{coordenadas de Jacobi}.

\subsection{Definición y motivación}
Normalmente describimos un sistema de \(N\) cuerpos con posiciones \(\vec{r}_i\) respecto a un origen arbitrario (6\(N\) coordenadas).  

En coordenadas de Jacobi se calculan centros de masa de pares sucesivos, de modo que el sistema se describe completamente usando \textbf{solo distancias relativas} (y los centros de masa de los subsistemas).

Para el caso de \textbf{dos cuerpos} (base del método):
\begin{align*}
\vec{R}_{\text{CM}} &= \frac{m_1 \vec{r}_1 + m_2 \vec{r}_2}{M} \\
\vec{r} &= \vec{r}_1 - \vec{r}_2
\end{align*}
donde \(M = m_1 + m_2\).

\subsection{Transformación de coordenadas}
\textbf{De coordenadas convencionales a Jacobi:}
\[
\vec{R}_{\text{CM}} = \frac{m_1 \vec{r}_1 + m_2 \vec{r}_2}{M}, \quad
\vec{r} = \vec{r}_1 - \vec{r}_2
\]

\textbf{De Jacobi a coordenadas convencionales:}
\[
\vec{r}_1 = \vec{R}_{\text{CM}} + \frac{m_2}{M} \vec{r}, \quad
\vec{r}_2 = \vec{R}_{\text{CM}} - \frac{m_1}{M} \vec{r}
\]

Estas transformaciones son idénticas a las derivadas en la clase anterior para el problema relativo, pero aquí se enfatiza su uso recursivo en sistemas jerárquicos de \(N > 2\).

\begin{importante}
Las coordenadas de Jacobi desacoplan el movimiento del centro de masa (rectilíneo uniforme) del movimiento relativo en cada nivel jerárquico. Cada vector relativo \(\vec{r}\) satisface la ecuación del problema de dos cuerpos:
\[
\ddot{\vec{r}} = -\frac{\mu}{r^3}\vec{r}, \quad \mu = G(m_1 + m_2)
\]
\end{importante}

\section{Ejemplo práctico: sistema jerárquico de tres cuerpos (Zul Aa + Ab + B)}
En el ejercicio de la clase se construye un sistema jerárquico triple:

\begin{itemize}[leftmargin=2em,label={}]
\item \textbf{Zul} (sistema total)
  \begin{itemize}[leftmargin=2em,label={}]
  \item \textbf{A} (subsistema interno)
    \begin{itemize}[leftmargin=2em,label={}]
    \item \textbf{Aa} (estrella primaria, \(m = 1.0\))
    \item \textbf{Ab} (planeta, \(m = 0.01\))
    \end{itemize}
  \item \textbf{B} (estrella secundaria, \(m = 5.0\))
  \end{itemize}
\end{itemize}

Se usan coordenadas de Jacobi descendiendo por el árbol:
\begin{enumerate}
    \item Fijar \(\vec{R}_{\text{CM}}\) y \(\vec{V}_{\text{CM}}\) del sistema total.
    \item Calcular posición y velocidad del CM del subsistema A respecto a B.
    \item Descender al interior de A: posición de Aa y Ab respecto al CM de A.
\end{enumerate}

Las ecuaciones exactas (implementadas en el cuaderno de clase) son:
\[
\vec{R}_{\text{CM,A}} = \vec{R}_{\text{CM,AB}} + \frac{m_B}{m_A + m_B} \vec{r}_{AB}
\]
\[
\vec{r}_B = \vec{R}_{\text{CM,AB}} - \frac{m_A}{m_A + m_B} \vec{r}_{AB}
\]
y luego dentro de A:
\[
\vec{r}_{\text{Aa}} = \vec{R}_{\text{CM,A}} + \frac{m_{\text{Ab}}}{m_A} \vec{r}_A, \quad
\vec{r}_{\text{Ab}} = \vec{R}_{\text{CM,A}} - \frac{m_{\text{Aa}}}{m_A} \vec{r}_A
\]

Este procedimiento se generaliza a cualquier \(N\).

\section{Consecuencias y perspectivas}
\begin{itemize}
    \item Los sistemas jerárquicos son la forma natural en que se presentan los sistemas estelares múltiples, planetarios y galácticos.
    \item Las coordenadas de Jacobi permiten inicializar simulaciones de \(N\) cuerpos de forma estable y física.
    \item En la parte 2 de la clase (próxima) se profundizará en la simulación numérica y estabilidad de estos sistemas.
\end{itemize}

\begin{importante}
La aproximación jerárquica es válida cuando las órbitas internas son mucho más rápidas y cerradas que las externas (condición de estabilidad jerárquica).
\end{importante}


\chapter{Clase 14: Sistemas Jerárquicos (parte 2)}

Esta clase es completamente práctica y se desarrolla íntegramente en el cuaderno de Jupyter \texttt{CuadernoClase\_SistemasJerarquicos\_Ejercicio(1).ipynb}. El objetivo es aplicar todo lo visto en la Clase 13: construir explícitamente un sistema jerárquico de tres cuerpos usando coordenadas de Jacobi y simular su evolución dinámica con la librería \texttt{pymcel}.

\section{Definición del sistema jerárquico de tres cuerpos}
Llamamos al sistema completo \textbf{Zul}. Sus componentes son:
\begin{itemize}
    \item \textbf{Zul Aa}: estrella primaria (\(m = 1.0\))
    \item \textbf{Zul Ab}: planeta (\(m = 0.01\))
    \item \textbf{Zul B}: estrella secundaria (\(m = 5.0\))
\end{itemize}

El árbol jerárquico es:

\begin{itemize}[leftmargin=2em,label={}]
\item \textbf{Zul} (sistema total)
  \begin{itemize}[leftmargin=2em,label={}]
  \item \textbf{A} (subsistema interno: Aa + Ab)
    \begin{itemize}[leftmargin=2em,label={}]
    \item \textbf{Aa} (estrella primaria, \(m_{\text{Aa}} = 1.0\))
    \item \textbf{Ab} (planeta, \(m_{\text{Ab}} = 0.01\))
    \end{itemize}
  \item \textbf{B} (estrella secundaria, \(m_{\text{B}} = 5.0\))
  \end{itemize}
\end{itemize}

\begin{importante}
Este es un sistema jerárquico típico: un planeta orbita una estrella (subsistema A) y ese subsistema orbita a su vez una estrella más masiva (B). Las masas y distancias están elegidas para que la jerarquía sea estable.
\end{importante}

\section{Condiciones iniciales en coordenadas de Jacobi}
Usamos unidades canónicas: \(G = 1\).

\subsection{Subsistema interno A (Aa + Ab)}
\[
m_{\text{A}} = m_{\text{Aa}} + m_{\text{Ab}} = 1.01
\]
\[
\vec{r}_{\text{A}} = \begin{pmatrix} 0.1 \\ 0 \\ 0 \end{pmatrix}, \quad
\vec{v}_{\text{A}} = \begin{pmatrix} 0 \\ 4.0 \\ 0 \end{pmatrix}
\]

\subsection{Vector relativo entre A y B}
\[
\vec{r}_{\text{AB}} = \begin{pmatrix} 2.0 \\ 0 \\ 0 \end{pmatrix}, \quad
\vec{v}_{\text{AB}} = \begin{pmatrix} 0 \\ 2.0 \\ 0 \end{pmatrix}
\]

\subsection{Centro de masa total del sistema}
\[
\vec{R}_{\text{CM}} = \begin{pmatrix} 0 \\ 0 \\ 0 \end{pmatrix}, \quad
\vec{V}_{\text{CM}} = \begin{pmatrix} 0.01 \\ 0 \\ 0 \end{pmatrix}
\]

\section{Descenso por el árbol jerárquico (cálculo paso a paso)}

\subsection{Nivel 1: Posiciones y velocidades del CM de A y de B}
\[
\vec{R}_{\text{CM,A}} = \vec{R}_{\text{CM}} + \frac{m_{\text{B}}}{m_{\text{A}} + m_{\text{B}}} \vec{r}_{\text{AB}} = \begin{pmatrix} 1.66389351 \\ 0 \\ 0 \end{pmatrix}
\]
\[
\vec{V}_{\text{CM,A}} = \vec{V}_{\text{CM}} + \frac{m_{\text{B}}}{m_{\text{A}} + m_{\text{B}}} \vec{v}_{\text{AB}} = \begin{pmatrix} 0.01 \\ 1.66389351 \\ 0 \end{pmatrix}
\]
\[
\vec{r}_{\text{B}} = \vec{R}_{\text{CM}} - \frac{m_{\text{A}}}{m_{\text{A}} + m_{\text{B}}} \vec{r}_{\text{AB}} = \begin{pmatrix} -0.33610649 \\ 0 \\ 0 \end{pmatrix}
\]
\[
\vec{v}_{\text{B}} = \vec{V}_{\text{CM}} - \frac{m_{\text{A}}}{m_{\text{A}} + m_{\text{B}}} \vec{v}_{\text{AB}} = \begin{pmatrix} 0.01 \\ -0.33610649 \\ 0 \end{pmatrix}
\]

\subsection{Nivel 2: Posiciones y velocidades de Aa y Ab dentro de A}
\[
\vec{r}_{\text{Aa}} = \vec{R}_{\text{CM,A}} + \frac{m_{\text{Ab}}}{m_{\text{A}}} \vec{r}_{\text{A}} = \begin{pmatrix} 1.66488361 \\ 0 \\ 0 \end{pmatrix}
\]
\[
\vec{v}_{\text{Aa}} = \vec{V}_{\text{CM,A}} + \frac{m_{\text{Ab}}}{m_{\text{A}}} \vec{v}_{\text{A}} = \begin{pmatrix} 0.01 \\ 1.70349747 \\ 0 \end{pmatrix}
\]
\[
\vec{r}_{\text{Ab}} = \vec{R}_{\text{CM,A}} - \frac{m_{\text{Aa}}}{m_{\text{A}}} \vec{r}_{\text{A}} = \begin{pmatrix} 1.56488361 \\ 0 \\ 0 \end{pmatrix}
\]
\[
\vec{v}_{\text{Ab}} = \vec{V}_{\text{CM,A}} - \frac{m_{\text{Aa}}}{m_{\text{A}}} \vec{v}_{\text{A}} = \begin{pmatrix} 0.01 \\ -2.29650253 \\ 0 \end{pmatrix}
\]

\begin{importante}
Estas fórmulas son las coordenadas de Jacobi aplicadas recursivamente. Cada nivel usa el centro de masa del subsistema anterior como nuevo origen.
\end{importante}

\section{Construcción del sistema para simulación}
El sistema completo se define como una lista de diccionarios:
\begin{lstlisting}
sistema = [
    dict(m = 1.0,   r = r_Zul_Aa, v = v_Zul_Aa),
    dict(m = 0.01,  r = r_Zul_Ab, v = v_Zul_Ab),
    dict(m = 5.0,   r = r_Zul_B,  v = v_Zul_B)
]
\end{lstlisting}

Se simula con la función \texttt{ncuerpos\_solución} de \texttt{pymcel} en el intervalo temporal:
\[
t \in [0, 50]\qquad \text{con 1000 puntos}
\]

\section{Resultados de la simulación y visualización}
La simulación muestra claramente el comportamiento jerárquico esperado:
\begin{itemize}
    \item \textbf{Órbita interna rápida}: el planeta Ab describe una órbita casi circular alrededor de Aa (radio \(\approx 0.1\)).
    \item \textbf{Órbita intermedia}: el subsistema A (Aa + Ab) describe una órbita alrededor del centro de masa total.
    \item \textbf{Órbita externa}: la estrella B orbita alrededor del centro de masa del sistema completo con un radio mayor (\(\approx 0.336\)).
\end{itemize}

Se generan gráficas de:
\begin{itemize}
    \item Trayectorias en el plano \(xy\) (todas las partículas).
    \item Trayectorias en el marco del centro de masa.
    \item Animación con \texttt{celluloid} que muestra el movimiento jerárquico completo.
\end{itemize}

\begin{importante}
La simulación numérica confirma que el sistema permanece estable durante el tiempo integrado. Esto valida la aproximación jerárquica y la correcta inicialización con coordenadas de Jacobi.
\end{importante}

\section{Generalización y consejos prácticos}
\begin{enumerate}
    \item Para cualquier sistema jerárquico de \(N\) cuerpos se sigue el mismo procedimiento recursivo: se desciende por el árbol calculando centros de masa y vectores relativos en cada nivel.
    \item Siempre se debe fijar primero el centro de masa y velocidad total del sistema (\(\vec{R}_{\text{CM}}\), \(\vec{V}_{\text{CM}}\)).
    \item La librería \texttt{pymcel} facilita enormemente la simulación una vez que se tienen las condiciones iniciales en el formato de lista de diccionarios.
    \item En sistemas reales (ej. sistemas estelares múltiples, exoplanetas circumbinarios) se usan exactamente estas mismas técnicas para generar condiciones iniciales estables.
\end{enumerate}

\section{Conclusiones de las Clases 13 y 14}
Las dos clases juntas muestran el camino completo:
\begin{itemize}
    \item Teoría (Clase 13): motivación, autoorganización y coordenadas de Jacobi.
    \item Práctica (Clase 14): construcción explícita de un sistema jerárquico y simulación numérica.
\end{itemize}

Esta metodología es la base para estudiar sistemas de \(N\) cuerpos más complejos en las próximas unidades.




\chapter{Clase 15: El problema relativo de dos cuerpos}

Esta clase, correspondiente a la Unidad 2 del curso (\emph{El problema de 2-Cuerpos}), se centra en la derivación rigurosa del \textbf{problema relativo} a partir del problema de \(N\) cuerpos. El objetivo es mostrar cómo, al considerar únicamente dos cuerpos, es posible reducir el sistema a una ecuación vectorial equivalente al movimiento de un solo cuerpo bajo una fuerza central, facilitando la solución analítica y numérica posterior.

\section{Motivación y contexto}
En el problema de \(N\) cuerpos (capítulo 6 del libro de texto), las ecuaciones de movimiento son acopladas y no tienen solución analítica general. Para \(N=2\), sin embargo, es posible transformar las ecuaciones a un sistema desacoplado: el movimiento del centro de masa (uniforme rectilíneo) y el movimiento \emph{relativo} entre los dos cuerpos. Este último es el que describe la órbita relativa y es equivalente al problema de Kepler.

El vector relativo se define como:
\[
\vec{r} = \vec{r}_1 - \vec{r}_2
\]
donde \(\vec{r}_1\) y \(\vec{r}_2\) son las posiciones de los cuerpos 1 y 2 en un marco inercial.

El centro de masa del sistema es:
\[
\vec{R}_{\text{CM}} = \frac{m_1 \vec{r}_1 + m_2 \vec{r}_2}{M}, \quad M = m_1 + m_2
\]

\section{De las posiciones a las velocidades}
Para expresar las posiciones individuales en términos del vector relativo y del centro de masa, resolvemos el sistema lineal:

\[
\vec{r}_1 - \vec{r}_2 = \vec{r}
\]
\[
m_1 \vec{r}_1 + m_2 \vec{r}_2 = M \vec{R}_{\text{CM}}
\]

La solución es:
\[
\vec{r}_1 = \vec{R}_{\text{CM}} + \frac{m_2}{M} \vec{r}
\]
\[
\vec{r}_2 = \vec{R}_{\text{CM}} - \frac{m_1}{M} \vec{r}
\]

Diferenciando con respecto al tiempo (velocidades):
\[
\dot{\vec{r}}_1 = \vec{V}_{\text{CM}} + \frac{m_2}{M} \dot{\vec{r}}
\]
\[
\dot{\vec{r}}_2 = \vec{V}_{\text{CM}} - \frac{m_1}{M} \dot{\vec{r}}
\]

\begin{importante}
Estas expresiones muestran que el movimiento de cada cuerpo es la superposición del movimiento uniforme del centro de masa más un término proporcional al movimiento relativo. El factor \(\frac{m_2}{M}\) indica que el cuerpo más ligero se mueve más en la coordenada relativa.
\end{importante}

\section{De las velocidades a las aceleraciones}
Diferenciando nuevamente obtenemos las aceleraciones:
\[
\ddot{\vec{r}}_1 = \frac{m_2}{M} \ddot{\vec{r}}
\]
\[
\ddot{\vec{r}}_2 = -\frac{m_1}{M} \ddot{\vec{r}}
\]

Estas relaciones son puramente cinemáticas y serán clave para relacionar las ecuaciones de movimiento newtonianas con la ecuación relativa.

\section{Ecuaciones de movimiento de las partículas (problema de \texorpdfstring{$N$}{N} cuerpos)}
Para dos cuerpos interactuando gravitacionalmente, las ecuaciones del problema de \(N\) cuerpos se reducen a:
\[
\ddot{\vec{r}}_i = -\sum_{j \neq i} \frac{G m_j}{r_{ij}^3} \vec{r}_{ij}
\]

Explicitamente:
\[
\ddot{\vec{r}}_1 = -\frac{G m_2}{r^3} (\vec{r}_1 - \vec{r}_2) = -\frac{\mu_2}{r^3} \vec{r}, \quad \mu_2 = G m_2
\]
\[
\ddot{\vec{r}}_2 = -\frac{G m_1}{r^3} (\vec{r}_2 - \vec{r}_1) = +\frac{\mu_1}{r^3} \vec{r}, \quad \mu_1 = G m_1
\]
donde \(r = |\vec{r}|\) y \(\vec{r} = \vec{r}_1 - \vec{r}_2\).

\section{Deducción de la ecuación de movimiento relativa}
Ahora combinamos las expresiones cinemáticas con las dinámicas. Hay dos caminos equivalentes; ambos se presentan para mayor claridad.

\textbf{Vía 1: Diferencia directa de aceleraciones}
\[
\ddot{\vec{r}} = \ddot{\vec{r}}_1 - \ddot{\vec{r}}_2 = \left( -\frac{\mu_2}{r^3} \vec{r} \right) - \left( +\frac{\mu_1}{r^3} \vec{r} \right) = -\frac{\mu_1 + \mu_2}{r^3} \vec{r}
\]
\[
\mu = G(m_1 + m_2) \implies \ddot{\vec{r}} = -\frac{\mu}{r^3} \vec{r}
\]

\textbf{Vía 2: Reemplazo y resta (como en las láminas de clase)}
Igualamos las dos expresiones para \(\ddot{\vec{r}}_1\):
\[
\frac{m_2}{M} \ddot{\vec{r}} = -\frac{\mu_2}{r^3} \vec{r}
\]
y para \(\ddot{\vec{r}}_2\):
\[
-\frac{m_1}{M} \ddot{\vec{r}} = +\frac{\mu_1}{r^3} \vec{r}
\]

Reordenando y restando (o resolviendo directamente) se obtiene el mismo resultado:
\[
\ddot{\vec{r}} = -\frac{\mu}{r^3} \vec{r}, \quad \mu = G(m_1 + m_2)
\]

Esta es la \textbf{ecuación del problema de los dos cuerpos relativo}.

\begin{importante}
La ecuación \(\ddot{\vec{r}} = -\dfrac{\mu}{r^3} \vec{r}\) describe el movimiento relativo como si fuera un cuerpo de masa reducida bajo una fuerza central inversa al cuadrado. Es idéntica a la ecuación de Kepler con parámetro gravitacional efectivo \(\mu = G(m_1 + m_2)\).
\end{importante}

\section{El problema de los dos cuerpos en coordenadas de Jacobi}
En el formalismo de Jacobi (útil para generalizaciones a \(N\) cuerpos), el vector relativo se expresa como:
\[
\vec{r} = \vec{r}_1 - \vec{r}_2
\]
y la ecuación de movimiento es exactamente la misma:
\[
\ddot{\vec{r}} = -\frac{\mu}{r^3} \vec{r}, \quad \mu \equiv G(m_1 + m_2)
\]

\textbf{Importante aclaración geométrica:}  
Al representar gráficamente la órbita relativa (la curva que traza el extremo del vector \(\vec{r}\)), el origen del sistema de coordenadas (el foco de la cónica) \textbf{no coincide con ningún objeto físico}. No es el centro de masa ni la posición del cuerpo más pesado. El punto \( \vec{r} = 0 \) corresponde al caso de colisión (ambos cuerpos en el mismo lugar), y el foco de la elipse/hipérbola es un punto vacío en el espacio físico. Esto evita confusiones comunes con la aproximación \(m_1 \gg m_2\), donde se fija el cuerpo masivo en el origen.

\begin{importante}
En coordenadas de Jacobi, el problema relativo es independiente del marco de referencia del centro de masa. El movimiento del CM es trivial (rectilíneo uniforme), mientras que toda la dinámica orbital reside en \(\vec{r}\).
\end{importante}

\section{Síntesis y consecuencias}
La derivación anterior muestra que:
\begin{itemize}
    \item El problema de dos cuerpos se reduce a una ecuación de segundo orden vectorial en 3D (6 variables de estado).
    \item El centro de masa se mueve con velocidad constante \(\vec{V}_{\text{CM}}\).
    \item Toda la información orbital (forma, tamaño, orientación) está contenida en la solución de \(\vec{r}(t)\).
\end{itemize}

Esta ecuación es la base para todos los desarrollos posteriores: ecuación de la trayectoria, anomalías, elementos orbitales, variables universales y perturbaciones (ver Clases 16 y siguientes).

\section{Ejemplo numérico ilustrativo}
Supongamos dos cuerpos con \(m_1 = 1\,M_\odot\), \(m_2 = 1\,M_\text{Tierra}\), separación inicial \(r = 1\,\text{UA}\), velocidad relativa perpendicular adecuada para órbita circular. La ecuación \(\ddot{\vec{r}} = -\mu/r^3 \vec{r}\) con \(\mu \approx G M_\odot\) (ya que \(m_2 \ll m_1\)) reproduce la órbita kepleriana con error negligible. Para masas comparables (ej. sistema binario), el factor completo \(\mu = G(m_1+m_2)\) es esencial.
\chapter{Clase 16: Cuadraturas del problema relativo de dos cuerpos}

Esta clase continúa el desarrollo analítico del problema relativo de dos cuerpos introducido en la clase anterior. El objetivo es identificar las \textbf{integrales primeras} o \textbf{cuadraturas} que permiten resolver completamente la dinámica relativa y establecer por qué este problema es integrable de forma exacta.

\section{Planteamiento del sistema y necesidad de cuadraturas}
Partimos de la ecuación diferencial que describe el movimiento relativo entre dos cuerpos sometidos a gravitación mutua:
\[
\ddot{\vec{r}} = -\frac{\mu}{r^3}\vec{r}
\]
donde \(\mu = G(m_1 + m_2)\) y \(r = |\vec{r}|\).

Esta es una ecuación diferencial vectorial de segundo orden. Al descomponerla en coordenadas cartesianas, el problema queda gobernado por \textbf{6 variables dinámicas independientes} \((x, y, z, \dot{x}, \dot{y}, \dot{z})\). Para integrar completamente el sistema y obtener una solución cerrada, es necesario hallar \textbf{6 cuadraturas} o integrales primeras independientes.

Aunque el problema de \(N\) cuerpos posee constantes de movimiento globales, el problema relativo de dos cuerpos constituye un sistema autónomo nuevo que requiere su propia deducción analítica paso a paso.

\begin{importante}
La meta de esta clase no es todavía escribir la órbita final en forma cónica, sino mostrar cómo la dinámica admite suficientes constantes de movimiento para cerrar completamente el sistema.
\end{importante}

\section{Primera cuadratura: momento angular relativo específico}
\subsection{Obtención mediante factor integrante}

Dado que el término de fuerza es proporcional a $\vec{r}$, el producto vectorial por la izquierda con $\vec{r}$ anula dicho término:

$$
\vec{r} \times \ddot{\vec{r}} = -\frac{\mu}{r^3}(\vec{r} \times \vec{r})
$$

Por la propiedad antisimétrica del producto vectorial, el producto de cualquier vector consigo mismo es nulo:
$$
\vec{r} \times \vec{r} = \vec{0} \quad \Rightarrow \quad \vec{r} \times \ddot{\vec{r}} = \vec{0}
$$

El siguiente paso consiste en reconocer que el lado izquierdo de esta igualdad es parte de la derivada temporal de un producto vectorial. Aplicamos la regla de Leibniz para la derivada de un producto cruz:
$$
\frac{d}{dt}(\vec{r} \times \dot{\vec{r}}) = \dot{\vec{r}} \times \dot{\vec{r}} + \vec{r} \times \ddot{\vec{r}}
$$
El primer término del lado derecho es idénticamente nulo porque el producto vectorial de un vector consigo mismo siempre se anula ($\dot{\vec{r}} \times \dot{\vec{r}} = \vec{0}$). Por lo tanto, la derivada se reduce a:
$$
\frac{d}{dt}(\vec{r} \times \dot{\vec{r}}) = \vec{r} \times \ddot{\vec{r}}
$$
Sustituyendo en la ecuación obtenida previamente:
$$
\frac{d}{dt}(\vec{r} \times \dot{\vec{r}}) = \vec{0}
$$

Esta igualdad indica que el vector $\vec{r} \times \dot{\vec{r}}$ tiene derivada temporal nula, lo cual implica que es constante en el tiempo. Integrando formalmente respecto al tiempo:
$$
\int \frac{d}{dt}(\vec{r} \times \dot{\vec{r}})\, dt = \int \vec{0}\, dt \quad \Rightarrow \quad \vec{r} \times \dot{\vec{r}} = \vec{C}
$$
donde $\vec{C}$ es un vector constante de integración. En mecánica celeste, a esta constante se le denomina \textbf{Momento Angular Relativo Específico} y se denota como:
$$
\boxed{\vec{h} \equiv \vec{r} \times \vec{v} = \text{cte.}}
$$
con $\vec{v} = \dot{\vec{r}}$.


Esta integral aporta \textbf{3 cuadraturas} (una por cada componente cartesiana). El vector \(\vec{h}\) es el \textbf{momento angular relativo específico}, es decir, por unidad de masa reducida.

\subsection{Interpretación física y geométrica}
\begin{itemize}
    \item \textbf{Movimiento plano:} como \(\vec{h}\) es constante y perpendicular tanto a \(\vec{r}\) como a \(\vec{v}\), el movimiento relativo queda confinado a un plano fijo perpendicular a \(\vec{h}\).
    \item \textbf{Coordenadas polares:} en el plano orbital, su módulo se expresa como \(h = r^2\dot{\theta}\).
    
    \textbf{Demostración}
    
    $$\vec{r} = r \vec{e}_r$$

$$\dot{\vec{r}} = \dot{r} \vec{e}_r + r\dot{\theta} \vec{e}_\theta$$

donde $\{\vec{e}_r, \vec{e}_\theta\}$ es la base ortonormal móvil del plano orbital. Calculando el momento angular relativo específico:

$$\vec{h} = \vec{r} \times \dot{\vec{r}} = (r \vec{e}_r) \times (\dot{r} \vec{e}_r + r\dot{\theta} \vec{e}_\theta)$$

Aplicando la distributividad del producto vectorial:

$$\vec{h} = r\dot{r}(\vec{e}_r \times \vec{e}_r) + r^2\dot{\theta}(\vec{e}_r \times \vec{e}_\theta)$$

Dado que $\vec{e}_r \times \vec{e}_r = \vec{0}$ y $\vec{e}_r \times \vec{e}_\theta = \vec{k}$ (vector unitario perpendicular al plano orbital, siguiendo la regla de la mano derecha):

$$\vec{h} = r^2\dot{\theta} \vec{k}$$

El módulo escalar de esta constante vectorial es directamente:

$$h = |\vec{h}| = r^2\dot{\theta}$$

    \item \textbf{Velocidad areal y segunda ley de Kepler:} el área infinitesimal barrida por el radio vector es \(dA = \frac{1}{2}r^2 d\theta\). Derivando respecto al tiempo:
\end{itemize}

\[
\frac{dA}{dt} = \frac{1}{2}r^2\dot{\theta} = \frac{h}{2} = \text{cte.}
\]

Esto demuestra que la velocidad areal es constante. La magnitud de \(h\) controla directamente la tasa de barrido del área, lo cual constituye la formulación dinámica de la segunda ley de Kepler.

\begin{importante}
La conservación del momento angular específico no solo reduce el número de grados efectivos del problema, sino que además garantiza que toda órbita kepleriana es plana.
\end{importante}

\section{Segunda cuadratura: energía relativa específica}
\subsection{Obtención mediante producto escalar}
Multiplicamos escalarmente la ecuación de movimiento por la velocidad \(\dot{\vec{r}}\):
\[
\dot{\vec{r}} \cdot \ddot{\vec{r}} = -\frac{\mu}{r^3}\vec{r} \cdot \dot{\vec{r}}
\]

\textbf{Deducción analítica del LHS (Término cinético)}

Partimos del producto escalar $\dot{\vec{r}} \cdot \ddot{\vec{r}}$. Recordemos que el cuadrado de la velocidad se define como el producto escalar de la velocidad consigo misma:
$$
v^2 \equiv \vec{v} \cdot \vec{v} = \dot{\vec{r}} \cdot \dot{\vec{r}}
$$
Diferenciamos esta expresión respecto al tiempo aplicando la regla de la derivada para un producto escalar:
$$
\frac{d}{dt}(v^2) = \frac{d}{dt}(\dot{\vec{r}} \cdot \dot{\vec{r}}) = \ddot{\vec{r}} \cdot \dot{\vec{r}} + \dot{\vec{r}} \cdot \ddot{\vec{r}}
$$
Dado que el producto escalar es conmutativo ($\vec{a} \cdot \vec{b} = \vec{b} \cdot \vec{a}$), ambos términos son idénticos:
$$
\frac{d}{dt}(v^2) = 2(\dot{\vec{r}} \cdot \ddot{\vec{r}})
$$
Despejando el producto original:
$$
\dot{\vec{r}} \cdot \ddot{\vec{r}} = \frac{1}{2}\frac{d}{dt}(v^2)
$$

\textbf{Deducción analítica del RHS (Término potencial)}

Partimos de la expresión $-\dfrac{\mu}{r^3}\vec{r} \cdot \dot{\vec{r}}$. Primero relacionamos el producto escalar $\vec{r} \cdot \dot{\vec{r}}$ con la derivada de la distancia $r$. Sabemos que:
$$
r^2 = \vec{r} \cdot \vec{r}
$$
Diferenciamos implícitamente respecto al tiempo:
$$
\frac{d}{dt}(r^2) = \frac{d}{dt}(\vec{r} \cdot \vec{r}) \quad \Rightarrow \quad 2r\dot{r} = 2\vec{r} \cdot \dot{\vec{r}}
$$
De aquí se obtiene la identidad fundamental:
$$
\vec{r} \cdot \dot{\vec{r}} = r\dot{r}
$$
Sustituimos esta identidad en el RHS:
$$
-\frac{\mu}{r^3}(\vec{r} \cdot \dot{\vec{r}}) = -\frac{\mu}{r^3}(r\dot{r}) = -\mu\frac{\dot{r}}{r^2}
$$
Ahora identificamos la derivada exacta de $1/r$. Aplicando la regla de la cadena:
$$
\frac{d}{dt}\left(\frac{1}{r}\right) = \frac{d}{dt}(r^{-1}) = -r^{-2}\frac{dr}{dt} = -\frac{\dot{r}}{r^2}
$$
Observamos que el término obtenido en el RHS coincide exactamente con esta derivada. Por lo tanto:
$$
-\mu\frac{\dot{r}}{r^2} = \mu\left(-\frac{\dot{r}}{r^2}\right) = \mu\frac{d}{dt}\left(\frac{1}{r}\right)
$$

\textit{Nota analítica:} El signo negativo presente en la ecuación original se absorbe en la derivada de $1/r$, dando como resultado un coeficiente positivo frente a la derivada exacta. Esto es crucial para obtener la expresión correcta de la energía)

Identificamos las derivadas exactas:
\[
\frac{1}{2}\frac{d}{dt}(v^2) = -\mu \frac{d}{dt}\left(\frac{1}{r}\right)
\]

Reordenando e integrando:
\[
\frac{d}{dt}\left(\frac{v^2}{2} - \frac{\mu}{r}\right) = 0
\]

Definimos la \textbf{energía relativa específica} como la constante resultante:
\[
\mathcal{E} \equiv \frac{v^2}{2} - \frac{\mu}{r} = \text{cte.}
\]

Esta integral escalar aporta \textbf{1 cuadratura} adicional.

\subsection{Interpretación física}
Representa la conservación de la energía mecánica por unidad de masa reducida. Su valor algebraico clasifica la trayectoria:
\begin{itemize}
    \item \(\mathcal{E} < 0\): órbita ligada (elipse).
    \item \(\mathcal{E} = 0\): caso límite de escape (parábola).
    \item \(\mathcal{E} > 0\): órbita no ligada (hipérbola).
\end{itemize}

\begin{importante}
El signo de \(\mathcal{E}\) permite distinguir si el movimiento permanece acotado o si el cuerpo escapa al infinito, incluso antes de resolver explícitamente la trayectoria.
\end{importante}

\section{Tercera cuadratura: vector de excentricidad}
\subsection{Construcción del factor integrante}
Para completar el sistema, cruzamos la ecuación de movimiento con el vector constante \(\vec{h}\):
\[
\ddot{\vec{r}} \times \vec{h} = -\frac{\mu}{r^3}\vec{r} \times (\vec{r} \times \vec{v})
\]

Aplicamos la identidad del triple producto vectorial (BAC-CAB):
\[
\vec{r} \times (\vec{r} \times \vec{v}) = \vec{r}(\vec{r} \cdot \vec{v}) - \vec{v}(\vec{r} \cdot \vec{r}) = r\dot{r}\vec{r} - r^2\vec{v}
\]

Sustituyendo y simplificando:
\[
\ddot{\vec{r}} \times \vec{h} = \mu\left(\frac{\vec{v}}{r} - \frac{\dot{r}\vec{r}}{r^2}\right)
\]

El término entre paréntesis coincide con la derivada temporal del vector unitario radial \(\hat{r} = \vec{r}/r\):
\[
\frac{d}{dt}\left(\frac{\vec{r}}{r}\right) = \frac{\dot{\vec{r}}}{r} - \frac{\vec{r}\dot{r}}{r^2}
\]

Por tanto, la ecuación se reduce a una derivada total exacta:
\[
\frac{d}{dt}(\dot{\vec{r}} \times \vec{h}) = \frac{d}{dt}\left(\mu\frac{\vec{r}}{r}\right)
\]

Integrando, obtenemos una nueva constante vectorial:
\[
\vec{e} \equiv \frac{\vec{v} \times \vec{h}}{\mu} - \frac{\vec{r}}{r} = \text{cte.}
\]

\subsection{Propiedades del vector de excentricidad}
Este vector, también llamado \textbf{vector de Laplace} o \textbf{vector de Runge--Lenz}, posee varias propiedades fundamentales:
\begin{itemize}
    \item yace en el plano orbital,
    \item es perpendicular a \(\vec{h}\),
    \item apunta hacia el periastro,
    \item su módulo coincide con la excentricidad orbital: \(e = |\vec{e}|\).
\end{itemize}

Formalmente aporta \textbf{3 cuadraturas}, aunque no todas son independientes de las integrales anteriores.

\begin{importante}
El vector de excentricidad contiene información geométrica fina sobre la orientación y la forma de la órbita. Mientras \(\vec{h}\) fija el plano orbital, \(\vec{e}\) fija la dirección del periastro dentro de ese plano.
\end{importante}

\section{Cierre del sistema: conteo de cuadraturas independientes}
Sumando las constantes obtenidas, tenemos:
\begin{itemize}
    \item momento angular relativo específico \(\vec{h}\): 3 componentes,
    \item energía relativa específica \(\mathcal{E}\): 1 escalar,
    \item vector de excentricidad \(\vec{e}\): 3 componentes.
\end{itemize}

El total aparente es de \(7\) ecuaciones. Sin embargo, el sistema no posee siete grados de libertad independientes, porque existen relaciones algebraicas que ligan estas constantes:
\begin{enumerate}
    \item ortogonalidad geométrica:
\end{enumerate}
\[
\vec{h} \cdot \vec{e} = 0
\]
\begin{enumerate}[resume]
    \item relación escalar derivada de \(\vec{e} \cdot \vec{e}\):
\end{enumerate}

$$e^2 = 1 + \frac{2\mathcal{E}h^2}{\mu^2}$$

\textbf{Demostración:}

Partimos de la definición del vector de excentricidad (vector de Laplace):
$$\vec{e} = \frac{\vec{v} \times \vec{h}}{\mu} - \frac{\vec{r}}{r}$$

Calculamos el producto escalar consigo mismo:
$$e^2 = \vec{e} \cdot \vec{e} = \left( \frac{\vec{v} \times \vec{h}}{\mu} - \frac{\vec{r}}{r} \right) \cdot \left( \frac{\vec{v} \times \vec{h}}{\mu} - \frac{\vec{r}}{r} \right)$$

Expandiendo algebraicamente:
$$e^2 = \frac{1}{\mu^2}(\vec{v} \times \vec{h}) \cdot (\vec{v} \times \vec{h}) - \frac{2}{\mu r}(\vec{v} \times \vec{h}) \cdot \vec{r} + \frac{\vec{r} \cdot \vec{r}}{r^2}$$

Evaluamos cada término por separado:
\begin{itemize}
	\item $\dfrac{\vec{r} \cdot \vec{r}}{r^2} = \dfrac{r^2}{r^2} = 1$.
	\item Usando la propiedad cíclica del triple producto escalar, $(\vec{v} \times \vec{h}) \cdot \vec{r} = \vec{h} \cdot (\vec{r} \times \vec{v})$. Por definición $\vec{h} = \vec{r} \times \vec{v}$, luego el producto es $\vec{h} \cdot \vec{h} = h^2$. El término completo queda $-\dfrac{2h^2}{\mu r}$.
	\item Aplicamos la identidad de Lagrange $|\vec{A} \times \vec{B}|^2 = A^2B^2 - (\vec{A}\cdot\vec{B})^2$:
   $$(\vec{v} \times \vec{h}) \cdot (\vec{v} \times \vec{h}) = v^2 h^2 - (\vec{v} \cdot \vec{h})^2$$
   Dado que $\vec{h} = \vec{r} \times \vec{v}$ es ortogonal a $\vec{v}$, se cumple $\vec{v} \cdot \vec{h} = 0$. Así, el término se reduce a $v^2 h^2$.
\end{itemize}

Sustituyendo los resultados en la expansión:
$$e^2 = \frac{v^2 h^2}{\mu^2} - \frac{2 h^2}{\mu r} + 1 = 1 + \frac{h^2}{\mu^2}\left( v^2 - \frac{2\mu}{r} \right)$$

Reconocemos la definición de Energía Relativa Específica $\mathcal{E} = \dfrac{v^2}{2} - \dfrac{\mu}{r}$. Multiplicando por 2 obtenemos $2\mathcal{E} = v^2 - \dfrac{2\mu}{r}$. Reemplazando en el paréntesis:

$$e^2 = 1 + \frac{2\mathcal{E}h^2}{\mu^2}$$

Esta expresión demuestra que la excentricidad geométrica es una función exclusiva de las integrales primeras $\mathcal{E}$ y $h$, cerrando analíticamente la relación entre la dinámica y la geometría de la cónica.
Al restar estas dependencias funcionales, el sistema queda con exactamente \textbf{6 cuadraturas independientes}, que coinciden con el número de variables dinámicas iniciales. Esto demuestra que el problema relativo de dos cuerpos es \textbf{completamente integrable} y permite derivar la ecuación de la trayectoria \(r(\theta)\) sin recurrir, en principio, a métodos numéricos.

\section{Síntesis}
En esta clase se obtuvo el conjunto fundamental de integrales primeras del problema relativo de dos cuerpos:
\begin{itemize}
    \item conservación del momento angular específico,
    \item conservación de la energía específica,
    \item conservación del vector de excentricidad.
\end{itemize}

Estas cuadraturas no solo permiten cerrar analíticamente el sistema, sino que preparan el camino para la deducción explícita de la ecuación de la cónica orbital y el estudio detallado de los elementos orbitales en la clase siguiente.
\chapter{Clase 17: Solución analítica del problema relativo y plano de Laplace}

Esta clase continúa el cierre analítico del problema relativo de dos cuerpos. A partir de las cuadraturas obtenidas en la clase anterior, se reduce el problema al plano orbital y se deriva la ecuación polar de la trayectoria. El objetivo central es mostrar que la solución relativa corresponde siempre a una sección cónica y que su geometría queda completamente determinada por las constantes del movimiento.

\section{Reducción al plano de la trayectoria}
Partimos del conteo de cuadraturas obtenido en el espacio tridimensional:
\begin{itemize}
    \item momento angular relativo específico \(\vec{h}\): 3 componentes,
    \item energía relativa específica \(\mathcal{E}\): 1 escalar,
    \item vector de Laplace o de excentricidad \(\vec{e}\): 3 componentes.
\end{itemize}

En apariencia esto produce 7 constantes, pero una de ellas queda ligada por la relación
\[
\vec{h} \cdot \vec{e} = 0,
\]
por lo que el sistema tiene finalmente \textbf{6 cuadraturas independientes}.

Al pasar al \textbf{plano de la trayectoria} —también llamado plano orbital o plano perpendicular a \(\vec{h}\)— el problema se reduce a dos dimensiones espaciales. En ese plano, el número de variables dinámicas independientes es \(4\): \((x, y, \dot{x}, \dot{y})\). En consecuencia, basta con describir el movimiento mediante cuatro constantes escalares equivalentes:
\begin{enumerate}
    \item \(h = |\vec{h}|\), módulo del momento angular específico,
    \item \(\mathcal{E}\), energía relativa específica,
    \item la dirección de \(\vec{e}\) en el plano orbital,
    \item el módulo \(e = |\vec{e}|\).
\end{enumerate}

Este marco permite enfocarse exclusivamente en la geometría y dinámica plana de la órbita.

\begin{importante}
La reducción al plano orbital no es una aproximación. Es una consecuencia exacta de la conservación del momento angular: toda la dinámica relativa ocurre en un único plano fijo.
\end{importante}

\section{Deducción de la ecuación polar de la trayectoria}
Proyectamos ahora el vector de Laplace sobre el vector posición \(\vec{r}\) para obtener la relación funcional entre \(r\) y \(\theta\).

\subsection{Proyección geométrica}
Por definición del producto punto:
\[
\vec{e} \cdot \vec{r} = e r \cos\theta
\]
donde \(\theta\) es la verdadera anomalía, es decir, el ángulo entre \(\vec{e}\) y \(\vec{r}\).

\subsection{Proyección algebraica}
Usando la definición del vector \(\vec{e}\):
\[
\vec{e} \cdot \vec{r} = \left(\frac{\vec{v} \times \vec{h}}{\mu} - \frac{\vec{r}}{r}\right) \cdot \vec{r}
\]
\[
\vec{e} \cdot \vec{r} = \frac{(\vec{v} \times \vec{h}) \cdot \vec{r}}{\mu} - \frac{\vec{r} \cdot \vec{r}}{r}
\]

Usando que \((\vec{v} \times \vec{h}) \cdot \vec{r} = h^2\) y que \(\vec{r} \cdot \vec{r} = r^2\), resulta:
\[
\vec{e} \cdot \vec{r} = \frac{h^2}{\mu} - r
\]

\subsection{Igualación y despeje de \texorpdfstring{$r(\theta)$}{r(theta)}}
Igualando la expresión geométrica con la algebraica:
\[
 e r \cos\theta = \frac{h^2}{\mu} - r
\]
\[
 r(1 + e\cos\theta) = \frac{h^2}{\mu}
\]

Definiendo el \textbf{parámetro} o \textbf{semilatus rectum}
\[
 p \equiv \frac{h^2}{\mu},
\]
se llega a la ecuación polar final de la trayectoria:
\[
\boxed{r(\theta) = \frac{p}{1 + e\cos\theta}}
\]

Esta es la ecuación general de una sección cónica con un foco en el origen del sistema relativo.

\begin{importante}
La ecuación \(r(\theta) = \frac{p}{1 + e\cos\theta}\) muestra que el problema relativo de dos cuerpos tiene como soluciones exactas cónicas: circunferencias, elipses, parábolas e hipérbolas aparecen como casos particulares según el valor de \(e\).
\end{importante}

\section{Resumen de la integración completa}
Los resultados principales de esta clase pueden sintetizarse así:
\begin{enumerate}
    \item El problema relativo de dos cuerpos posee suficientes cuadraturas independientes para ser completamente integrable.
    \item En el plano orbital, la dinámica queda caracterizada por \(h\), \(\mathcal{E}\), la dirección de \(\vec{e}\) y su magnitud \(e\).
    \item El vector \(\vec{e}\) fija la orientación de la órbita y su módulo determina la excentricidad geométrica mediante
\[
 e^2 = 1 + \frac{2\mathcal{E}h^2}{\mu^2}.
\]
    \item La proyección de \(\vec{e}\) sobre \(\vec{r}\) conduce directamente a
\[
 r(\theta) = \frac{p}{1 + e\cos\theta},
\]
lo que demuestra analíticamente que la trayectoria relativa es siempre una cónica.
\end{enumerate}

Este desarrollo cierra la solución analítica del problema de dos cuerpos y deja listo el marco conceptual para estudiar con más detalle la geometría orbital, los elementos de la cónica y su representación en distintos sistemas de referencia.
\chapter{Clase 18: Geometría de las cónicas (parte 1)}

Esta clase inicia el estudio geométrico de las cónicas desde la perspectiva de la mecánica celeste. Después de haber deducido la ecuación polar de la trayectoria relativa, el siguiente paso es comprender en qué sistema de referencia esa ecuación adopta su forma más simple y cómo se transforma al sistema de referencia del observador. El objetivo es establecer con claridad el \textbf{sistema natural de la cónica} y desarrollar la herramienta algebraica básica para relacionarlo con un sistema externo: la \textbf{rotación plana}.

\section{Introducción y repaso del marco teórico}
El problema relativo de dos cuerpos, una vez integradas sus cuadraturas, conduce a la ecuación de la trayectoria en coordenadas polares:
\[
 r(\theta) = \frac{p}{1 + e\cos\theta}
\]
donde \(p = h^2/\mu\) es el parámetro semilatus rectum, \(e = |\vec{e}|\) es la excentricidad y \(\theta\) es la verdadera anomalía medida desde la dirección del vector de Laplace \(\vec{e}\).

Esta expresión describe completamente la forma de la órbita, pero está escrita en un sistema de coordenadas muy particular: aquel cuyo eje polar coincide con la dirección de \(\vec{e}\) y cuyo origen se sitúa en el foco gravitacional. A este marco se le denomina \textbf{sistema natural de la cónica}.

El objetivo analítico de esta clase es formalizar la geometría de las secciones cónicas en su sistema natural, identificar las limitaciones que presenta para la descripción observacional y construir rigurosamente las herramientas de rotación plana necesarias para pasar del sistema natural al sistema de referencia del observador.

\begin{importante}
La ecuación polar de la cónica no depende solo de la forma geométrica de la órbita, sino también de la elección del sistema de referencia. Su forma simple se obtiene únicamente en el sistema natural alineado con la órbita.
\end{importante}

\section{El sistema natural de la cónica}
Toda sección cónica posee una geometría intrínseca que se describe de manera más simple cuando se adopta un sistema de ejes alineado con sus elementos de simetría.

\subsection{Definición del marco natural \texorpdfstring{$(x',y')$}{(x',y')}}
Se define un sistema cartesiano plano \((x', y')\) con las siguientes características:
\begin{itemize}
    \item el eje \(x'\) coincide con el \textbf{eje de simetría} de la cónica,
    \item el origen puede ubicarse en el \textbf{foco} (sistema focal) o en el \textbf{centro geométrico} (sistema central), dependiendo de la conveniencia analítica,
    \item el eje \(y'\) es perpendicular a \(x'\) y completa la base ortonormal derecha del plano orbital.
\end{itemize}

En este sistema, la descripción analítica se simplifica notablemente porque la cónica presenta simetría respecto a \(x'\) y \(y'\).

\subsection{Ecuaciones en el sistema natural}
Dos descripciones típicas en este sistema son:
\begin{itemize}
    \item \textbf{Ecuación polar con origen en el foco:}
\[
 r = \frac{p}{1 + e\cos\theta'}
\]
    donde \(\theta'\) es el ángulo medido desde el semieje positivo \(x'\) hasta el vector posición \(\vec{r}\).
    \item \textbf{Ecuaciones paramétricas cartesianas:} dependiendo del tipo de cónica, las coordenadas \((x', y')\) pueden expresarse en función de un parámetro auxiliar. Por ejemplo, para la elipse usando la anomalía excéntrica \(E\):
\[
 x' = a(\cos E - e), \quad y' = a\sqrt{1-e^2}\sin E
\]
\end{itemize}

Estas expresiones conservan su forma funcional dentro del sistema natural, independientemente de cómo esté orientada la órbita en el espacio físico.

\section{El sistema de referencia del observador y el problema de alineación}
En la práctica astronómica, las posiciones y velocidades se miden respecto a un sistema de referencia fijo o inercial, denominado \textbf{sistema del observador} $(x,y,z)$. Este sistema puede ser, por ejemplo:
\begin{itemize}
    \item el plano ecuatorial terrestre,
    \item el plano de la eclíptica,
    \item el plano tangente a la esfera celeste en la dirección de un objeto.
\end{itemize}

El problema fundamental es que casi nunca el sistema de ejes utilizado en las observaciones coincide con el sistema natural de la órbita. El plano orbital puede estar inclinado respecto al plano de referencia y la línea de ápsides, es decir, la dirección de $\vec{e}$, puede formar un ángulo arbitrario con el eje $x$ del observador.

Por lo tanto, para comparar predicciones teóricas con datos observacionales, es necesario establecer una transformación matemática rigurosa que relacione las coordenadas del sistema natural $(x',y',z')$ con las del sistema del observador $(x,y,z)$. Como el movimiento orbital es estrictamente plano, la transformación se reduce inicialmente a una \textbf{rotación en el plano orbital}, la cual se embebe naturalmente en el espacio tridimensional mediante una rotación alrededor del eje $z$ perpendicular al plano.

\begin{importante}
La dificultad observacional no está en la ecuación de la órbita, sino en que los datos rara vez están expresados en el sistema donde esa ecuación adopta su forma más simple.
\end{importante}

\section{Rotación de la cónica en el espacio tridimensional}
Consideremos dos sistemas cartesianos ortonormales que comparten el mismo origen:
\begin{itemize}
    \item sistema del observador: base $\{\hat{i}, \hat{j}, \hat{k}\}$,
    \item sistema natural de la cónica: base $\{\hat{i}', \hat{j}', \hat{k}'\}$.
\end{itemize}

El sistema natural está rotado un ángulo $\theta$ respecto al sistema del observador, alrededor del eje $z$, en sentido antihorario. El vector posición $\vec{r}$ es el mismo objeto geométrico, pero sus componentes cambian con la base elegida:
$$ \vec{r} = x\hat{i} + y\hat{j} + z\hat{k} = x'\hat{i}' + y'\hat{j}' + z'\hat{k}' $$

\subsection{Relación entre vectores unitarios}
Para expresar los componentes en una base en función de los de la otra, proyectamos los vectores unitarios rotados sobre la base fija. Dado que la rotación ocurre en el plano $x$-$y$, el eje $z$ permanece invariante:
$$ \hat{i}' = \cos\theta\,\hat{i} + \sin\theta\,\hat{j} $$
$$ \hat{j}' = -\sin\theta\,\hat{i} + \cos\theta\,\hat{j} $$
$$ \hat{k}' = \hat{k} $$

Sustituyendo en $\vec{r} = x'\hat{i}' + y'\hat{j}' + z'\hat{k}'$ y agrupando términos respecto a $\hat{i}$, $\hat{j}$ y $\hat{k}$, obtenemos:
$$ x\hat{i} + y\hat{j} + z\hat{k} = (x'\cos\theta - y'\sin\theta)\hat{i} + (x'\sin\theta + y'\cos\theta)\hat{j} + z'\hat{k} $$

Igualando componentes, se deduce la \textbf{regla de rotación directa en 3D}:
$$ \begin{cases} 
 x = x'\cos\theta - y'\sin\theta \\ 
 y = x'\sin\theta + y'\cos\theta \\ 
 z = z' 
\end{cases} $$

\section{Derivación analítica de la matriz de rotación \texorpdfstring{$\mathbf{R}_z(\theta)$}{Rz(theta)}}
El sistema anterior puede escribirse de manera compacta mediante álgebra matricial. Definimos los vectores columna de coordenadas tridimensionales:
$$ \begin{pmatrix} x \\ y \\ z \end{pmatrix} = \mathbf{R}_z(\theta)\begin{pmatrix} x' \\ y' \\ z' \end{pmatrix} $$

Aquí, $\mathbf{R}_z(\theta)$ es la \textbf{matriz de rotación en el espacio tridimensional alrededor del eje $z$}:
$$ \mathbf{R}_z(\theta) = \begin{pmatrix} 
\cos\theta & -\sin\theta & 0 \\ 
\sin\theta & \cos\theta & 0 \\ 
0 & 0 & 1 
\end{pmatrix} $$

\subsection{Propiedades analíticas de \texorpdfstring{$\mathbf{R}_z(\theta)$}{Rz(theta)}}
Esta matriz posee propiedades fundamentales en $\mathbb{R}^3$:
\begin{enumerate}
    \item \textbf{Ortogonalidad:}
$$ \mathbf{R}_z^T(\theta)\mathbf{R}_z(\theta) = \mathbf{I}_3 $$
    lo cual garantiza que la transformación preserva productos escalares, distancias y ángulos en el espacio tridimensional.
    \item \textbf{Determinante unitario:}
$$ \det(\mathbf{R}_z) = \cos^2\theta + \sin^2\theta = 1 $$
    indicando que la rotación es propia y no incluye reflexiones ni inversiones de paridad.
    \item \textbf{Preservación de la norma:}
$$ x^2 + y^2 + z^2 = x'^2 + y'^2 + z'^2 $$
    consistente con la invariancia de la distancia radial $r = |\vec{r}|$ y la estructura métrica euclidiana.
\end{enumerate}

Esta matriz es la herramienta fundamental para trasladar cualquier descripción geométrica de la cónica desde su sistema natural al sistema del observador, manteniendo la coherencia tridimensional.

\section{La rotación inversa y el cambio de signo angular}
Con frecuencia necesitamos hacer el camino contrario: dada una posición observada $(x,y,z)$, queremos recuperar sus coordenadas en el sistema natural $(x',y',z')$ para aplicar directamente la ecuación polar o las leyes de Kepler.

La transformación inversa se obtiene multiplicando por la inversa de la matriz de rotación:
$$ \begin{pmatrix} x' \\ y' \\ z' \end{pmatrix} = \mathbf{R}_z^{-1}(\theta)\begin{pmatrix} x \\ y \\ z \end{pmatrix} $$

Como $\mathbf{R}_z(\theta)$ es ortogonal, su inversa coincide con su traspuesta:
$$ \mathbf{R}_z^{-1}(\theta) = \mathbf{R}_z^T(\theta) = \begin{pmatrix} 
\cos\theta & \sin\theta & 0 \\ 
-\sin\theta & \cos\theta & 0 \\ 
0 & 0 & 1 
\end{pmatrix} $$

Analíticamente, esto equivale a cambiar el signo del ángulo:
$$ \mathbf{R}_z^{-1}(\theta) = \mathbf{R}_z(-\theta) $$

Por tanto, la rotación inversa en 3D se escribe como:
$$ \begin{cases} 
 x' = x\cos\theta + y\sin\theta \\ 
 y' = -x\sin\theta + y\cos\theta \\ 
 z' = z 
\end{cases} $$

\subsection{Unicidad de la matriz de transformación}
Solo existe una matriz de rotación plana que conecta ambos marcos. La misma matriz \(\mathbf{R}_z(\theta)\), o su traspuesta según la dirección del cambio de base, se utiliza tanto para pasar del sistema del observador al sistema natural como para hacer la transformación inversa.

Esta propiedad simplifica el tratamiento analítico y computacional, ya que no se requieren transformaciones distintas para ida y vuelta: basta con ajustar el signo del ángulo o trasponer la matriz.

\section{Síntesis}
Los resultados centrales de esta clase son:
\begin{enumerate}
    \item La ecuación polar \(r(\theta) = p/(1+e\cos\theta)\) y las ecuaciones paramétricas son válidas de forma natural en el sistema alineado con la cónica.
    \item Las observaciones astronómicas se realizan en un sistema de referencia que, en general, no coincide con el sistema natural de la órbita.
    \item La relación entre ambos sistemas se resuelve mediante una rotación plana descrita por la matriz \(\mathbf{R}_z(\theta)\), que es ortogonal, de determinante unitario y preserva la métrica euclidiana.
    \item La transformación inversa se obtiene con \(\theta \to -\theta\), garantizando la reversibilidad analítica del cambio de coordenadas.
\end{enumerate}

Este marco bidimensional es el primer paso hacia la descripción completa de la órbita en el espacio tridimensional. En clases posteriores, esta rotación plana se combinará con rotaciones adicionales asociadas al nodo ascendente y a la inclinación orbital para construir la matriz completa que conecta el sistema perifocal con un sistema inercial de referencia.

\begin{importante}
La geometría de las cónicas no solo determina la forma de la trayectoria orbital. También fija la estructura algebraica de las transformaciones de coordenadas que hacen posible comparar teoría, simulación y observación.
\end{importante}
\input{clases/clase20}

% ---------- FINAL ----------
\backmatter
% Puedes agregar bibliografía o anexos aquí si deseas.

\end{document}